\section{Short stirring bursts on the torus}\label{sec:compact}

This section separates the diffusion argument from the construction of
a planar machine. The input is a smooth periodic planar field and a
distinguished initial point. Its orbit either enters a target at an
integer phase or stays a positive distance away at every phase. The
output is a force for which the third velocity component distinguishes
these two cases.

We first work at viscosity one. Let $V(s,Y)$ be an effectively specified
smooth divergence-free field on $\R\times\T^2$, periodic of period one
and zero on a neighborhood of every integer phase. Its flow from phase
zero to phase $s$ is denoted $\Psi_s$. Assume that effective uniform
derivative bounds are available; choose an integer $L\ge1$ such that
\begin{equation}\label{eq:processor-derivatives}
 \|D_YV\|_{\infty,\mathrm{op}},\quad
 \|D_Y^2V\|_{\infty,\mathrm{op}}\le L,
 \qquad R=4^L.
\end{equation}
All derivative norms in this section are Euclidean operator norms.
Let $p\in\T^2$ be computable and let $E\subset\T^2$ be an open target.

\begin{theorem}[Injection, stirring and waiting]\label{thm:compact-detector}
For $V,p,E$ as above, choose either of the two pairs
\begin{equation}\label{eq:two-constants}
 (d,H)=(1/32,10^9)\quad\hbox{or}\quad(1/16,10^8).
\end{equation}
There is a smooth force on $\T^3$, effectively determined by $V,p$,
for which zero initial
velocity has a global smooth solution, unique in the classical torus
class of Section~\ref{sec:introduction}, with pressure zero and
the following properties.
\begin{enumerate}
\item If $\Psi_k(p)\in E$ for some integer $k\ge0$, then
$u_3(t,Y,z)>1/2$ for some finite $t$ and some $Y\in E$.
\item If $\dist_{\T^2}(\Psi_s(p),E)\ge2d$ for every $s\ge0$,
then $u_3(t,Y,z)<1/2$ for every $t\ge0$ and $Y\in E$.
\end{enumerate}
The third component is nonnegative. At viscosity one, its total
horizontal mass is less than $10^{-11}$. The same conclusions and
threshold hold for every positive computable viscosity, with the time
and horizontal velocity rescaled as in~\eqref{eq:compact-viscosity}.
The force is supported in the union of the horizontal support of $V$
and one fixed small neighborhood of $p$, times $\T$. All smoothness and
uniqueness assertions concern every finite time interval.
\end{theorem}

The force does not depend on $E$, so no effective presentation of the
open target is required. The target enters only the two orbit tests.
Here $d$ is the clearance required between a transported bump and the
target, and $H$ is an upper bound for the heat kernel at that clearance
or after one unit of waiting. The proof establishes that bound in
Lemma~\ref{lem:heat-bounds}.
The nonhalting hypothesis concerns the entire orbit. Integer-time
separation alone would leave open a false detection during a transition.
The first numerical pair serves the fixed torus detector in
Section~\ref{sec:compact-application}; the second serves the
slow-clock application identified at the end of that section.
The scalar proof is the same for both.

\subsection{An injection and a finite computation in each block}
\label{sec:compact-blocks}

Use the effective nondecreasing smooth step
\[
 b(r)=\begin{cases}e^{-1/r},&r>0,\\0,&r\le0,\end{cases}
 \qquad \theta(r)=\frac{b(r)}{b(r)+b(1-r)}.
\]
It is zero for $r\le0$ and one for $r\ge1$, with all derivatives flat
at the endpoints. Define a fixed bump on $\R^2$ by
\[
 g(Z)=\prod_{j=1}^2\theta(2(Z_j+1))\theta(2(1-Z_j)).
\]
It lies in $[0,1]$, equals one near zero, and is supported in $[-1,1]^2$.
Let an effective integer $C\ge1$ bound its first and second derivative
norms. For $n\ge1$, put
\begin{equation}\label{eq:compact-scales}
 \begin{aligned}
 b_n&=10^{-6}(2R)^{-n},&g_n(Y)&=g((Y-p)/b_n),\\
 D_n&=2Cb_n^{-2}(1+nL)R^{3n},&
 \delta_n&=\frac1{100(1+D_n)},\qquad t_n=2n.
 \end{aligned}
\end{equation}
The width $b_n$ makes the bump's mass small and its transported support
narrow. The quantity $D_n$ is chosen to bound the Laplacian of that
transported bump; Section~\ref{sec:compact-error} verifies the bound.
The duration $\delta_n$ then makes the accumulated diffusion error
at most $2\delta_nD_n<1/16$. These are separate roles: spatial width
controls the mass and separation, while physical duration controls
the error in height.
Here $g_n$ is periodized from the small square about a lift of $p$;
the copies have disjoint supports. In particular
\begin{equation}\label{eq:bump-bounds}
 \|\nabla g_n\|_\infty\le Cb_n^{-1},\qquad
 \|D^2g_n\|_\infty\le Cb_n^{-2},\qquad
 \int_{\T^2}g_n\le4b_n^2.
\end{equation}

Block $n$ injects this bump over time $\delta_n$, runs $n$ complete
periods of $V$ over another time $\delta_n$, and then waits for the
next block. In particular, its heat-only wait has length
$2-2\delta_n>1$. The time diagram in Figure~\ref{fig:blocks} shows the
roles of these intervals; their plotted lengths are schematic.

\begin{figure}[ht]
\centering
\begin{tikzpicture}[x=1cm,y=1cm,>=Latex,font=\small]
\draw[->] (0,0)--(12,0) node[right] {$t$};
\draw[fill=blue!12] (0.4,.15) rectangle (2.5,1.0);
\draw[fill=green!12] (2.5,.15) rectangle (5.2,1.0);
\draw[fill=gray!12] (5.2,.15) rectangle (11.5,1.0);
\node at (1.45,.58) {inject $g_n$};
\node[align=center,text width=2.4cm] at (3.85,.58) {$n$ processor\\periods};
\node at (8.35,.58) {heat evolution};
\node[below] at (.4,0) {$t_n$};
\node[below] at (2.5,0) {$t_n+\delta_n$};
\node[below] at (5.2,0) {$t_n+2\delta_n$};
\node[below] at (11.5,0) {$t_{n+1}$};
\node[above] at (1.45,1.03) {height one};
\node[above,align=center,text width=2.4cm] at (3.85,1.03) {small diffusion\\error};
\node[above,align=center,text width=4.8cm] at (8.35,1.03) {old mass becomes uniformly small};
\end{tikzpicture}
\caption{One torus block. Fresh scalar is injected at the initial code,
then transported through a longer finite prefix. The wait exceeds one
unit of heat time. The two active intervals are much shorter than the
wait; their relative sizes in the diagram are not to scale.}
\label{fig:blocks}
\end{figure}

Let $V_0$ be the single period of $V$ extended by zero outside $[0,1]$.
The phase collars make this extension smooth. Prescribe
\begin{equation}\label{eq:compact-program}
 \begin{split}
 s_n(t)&=n(t-t_n-\delta_n)/\delta_n,\\
 a(t,Y)&=\sum_{n\ge1}\frac n{\delta_n}
                  \sum_{j=0}^{n-1}V_0(s_n(t)-j,Y),\\
 h(t,Y)&=\sum_{n\ge1}\delta_n^{-1}
           \theta'((t-t_n)/\delta_n)g_n(Y),\\
 F(t,Y)&=\partial_ta+(a\cdot\nabla_Y)a-\Delta_Ya,
 \qquad f(t,Y,z)=(F(t,Y),h(t,Y)).
 \end{split}
\end{equation}
Each summand is supported in its indicated block. The sums are locally
finite in physical time and join smoothly at all endpoints.
Moreover $\diver_Ya=0$, $h\ge0$, and $a=h=0$ near time zero.
This prescription uses $V$ and the injection profiles, without using
the scalar that it will produce.

\subsection{Constructing the scalar and controlling old mass}

\begin{lemma}[Scalar evolution on the torus]\label{lem:torus-scalar}
For smooth $a,h$ on $[0,T]\times\T^2$ with $\diver_Ya=0$, smooth
initial data give a unique smooth solution of
$w_t+a\cdot\nabla_Yw=\Delta_Yw+h$.
Nonnegative initial data and source give $w\ge0$.
Homogeneous evolution contracts the supremum norm. On an interval
starting at $t_0$, zero initial data and $|h|\le A$ imply
$\|w(t)\|_\infty\le A(t-t_0)$.
\end{lemma}
\begin{proof}
Let $P_J$ project Fourier series onto $|k|_\infty\le J$ and solve
the finite linear system
\[
 \partial_tw_J=\Delta w_J+P_J(-a\cdot\nabla w_J+h),
 \qquad w_J(0)=P_Jw(0).
\]
For each integer $m\ge0$, differentiate through order $m$ and take
$L^2$ inner products with the corresponding derivatives of $w_J$.
The projection commutes with differentiation and drops out of these
inner products. The leading transport term has zero integral because
$\diver a=0$. In each remaining commutator a positive derivative
falls on $a$, leaving at most $m$ derivatives on $w_J$. Cauchy--Schwarz,
the finite-time bounds on $a,h$, and the nonpositive diffusion term give
\[
 \frac d{dt}\|w_J\|_{H^m}^2
       \le C_{m,T}(1+\|w_J\|_{H^m}^2).
\]
Gronwall bounds every $H^m$ norm uniformly in $J$ on $[0,T]$.

The equations bound the time derivative of every fixed Fourier
coefficient. A diagonal subsequence therefore converges uniformly in
time coefficientwise. The $H^{\ell+2}$ bound controls the differentiated
Fourier tail by a constant times
\[
 \left(\sum_{|k|_\infty>K}(1+|k|^2)^{-2}\right)^{1/2},
\]
which tends to zero in two dimensions. Hence convergence holds in
$C([0,T];C^\ell)$ for every $\ell$. The same tail argument applies
to the products and projections in the time-integrated equation.
Passing to that equation gives the scalar PDE with all spatial
derivatives. It then gives one time derivative, and repeated time
differentiation gives smoothness up to the initial time.

For comparison, subtract a proposed constant or linear upper bound
and then $\epsilon(t-t_0)$, initially allowing an arbitrarily small
strict gap as well. A first positive maximum has zero gradient,
nonpositive Laplacian and nonnegative time derivative, contradicting
the strict differential inequality. Let the gaps tend to zero. Apply
this argument also to $-w$ and to a difference of solutions. It proves
the bounds, positivity and uniqueness.
\end{proof}

Apply the lemma to~\eqref{eq:compact-program} with $w(0)=0$.
Solutions on finite intervals agree by uniqueness, giving a global
smooth nonnegative scalar. The velocity $(a,w)$ with pressure zero
is the unique three-dimensional solution by
Section~\ref{sec:scalar-structure}. Integrating the scalar equation
and using~\eqref{eq:bump-bounds} gives
\begin{equation}\label{eq:small-mass}
 0\le\int_{\T^2}w(t,Y)\dd Y
 \le M_0:=\sum_{n\ge1}4b_n^2
 =\frac{4\cdot10^{-12}}{4R^2-1}<10^{-11}.
\end{equation}
Thus infinitely many height-one injections still have very small total
mass. The heat kernel converts this mass bound into control of their
older contributions.

\begin{lemma}[Two explicit heat bounds]\label{lem:heat-bounds}
For either $(d,H)$ in~\eqref{eq:two-constants}, the unit two-torus
heat kernel satisfies
\begin{equation}\label{eq:heat-bound}
 K(\tau,Z)\le H
 \quad\hbox{if }\tau\ge1
 \quad\hbox{or}\quad\dist_{\T^2}(Z,0)\ge d.
\end{equation}
\end{lemma}
\begin{proof}
Periodizing the planar Gaussian gives
\[
 K(\tau,Z)=\frac1{4\pi\tau}
       \sum_{k\in\Z^2}e^{-|Z+k|^2/(4\tau)}.
\]
Its convolution solves the heat equation and tends uniformly to smooth
initial data as $\tau\downarrow0$, so uniqueness identifies it as the
heat kernel. Choose $Z\in[-1/2,1/2]^2$. The sup-norm shell
$|k|_\infty=j$ has $8j$ points and $|Z+k|\ge j/2$ for $j\ge1$.
Consequently
\[
 \sum_{k\in\Z^2}e^{-|Z+k|^2/8}
 \le1+8\sum_{j\ge1}j e^{-j/32}<10^4.
\]
For the last inequality, $e^{-1/32}\le32/33$ and
$\sum j q^j=q/(1-q)^2$ give an upper bound $1+8\cdot1056<10^4$.
For $0<\tau\le1$ and distance at least $d$, split each Gaussian
factor in half. One half satisfies
$\tau^{-1}e^{-d^2/(8\tau)}\le8/d^2$; summing the other halves
uses the preceding estimate. Thus
\[
 K(\tau,Z)\le\frac{8\cdot10^4}{4\pi d^2},
\]
which is less than $10^9$ for $d=1/32$ and less than $10^8$ for
$d=1/16$. At time one the same shell estimate applies without a
distance restriction; the maximum principle extends its upper bound
to all later times.
\end{proof}

Before block $n>1$, there has been more than one unit of heat-only
evolution. Therefore
\begin{equation}\label{eq:old-mass}
 0\le w(t_n,\cdot)\le HM_0;
\end{equation}
for $n=1$ the scalar is zero. From $t_n$ to $t_{n+1}$, linearity
splits $w=w^{\mathrm{old}}+w^{\mathrm{new}}$. The old part has
initial value $w(t_n)$ and no source, so remains between zero and
$HM_0$. The new part starts at zero and receives just the current
injection. It remains to estimate this new part near its transported
bump and away from it.

\subsection{A controlled diffusion error during the burst}\label{sec:compact-error}

The variational equations for $\Psi_s$ give, for $0\le s\le n$,
\begin{equation}\label{eq:flow-bounds}
 \|D\Psi_s^{\pm1}\|_\infty\le e^{Ln}\le R^n,
 \qquad
 \|D^2\Psi_s^{\pm1}\|_\infty
       \le nLe^{3Ln}\le nLR^{3n}.
\end{equation}
For the first bound, differentiate the flow once and use Gronwall.
The second variation has zero initial data, linear coefficient at
most $L$, and source at most $L$ times the square of the first
variation. Variation of constants bounds it by
$\int_0^s e^{L(s-r)}Le^{2Lr}\,dr$, and hence by $nLe^{3Ln}$.
The inverse is the flow of $-V(s-r,\cdot)$ over the reversed interval,
which obeys the same derivative bounds. These estimates are uniform
over all initial points; they do not require following the selected
machine execution.

During the two active intervals of block $n$, define the scalar that
would evolve without diffusion:
\begin{equation}\label{eq:inviscid-reference}
 w^0(t,Y)=
 \begin{cases}
 \theta((t-t_n)/\delta_n)g_n(Y),
                 &t_n\le t\le t_n+\delta_n,\\
 g_n(\Psi_{s_n(t)}^{-1}(Y)),
                 &t_n+\delta_n\le t\le t_n+2\delta_n.
 \end{cases}
\end{equation}
The endpoint flatness and phase collars make these expressions agree
smoothly. They satisfy $w^0_t+a\cdot\nabla w^0=h$.
The chain rule and~\eqref{eq:flow-bounds} bound the Hessian during
stirring by
\[
 Cb_n^{-2}R^{2n}+Cb_n^{-1}nLR^{3n}.
\]
Taking the trace costs at most two. Since $b_n<1$ and $R\ge1$,
the result is at most $D_n$ from~\eqref{eq:compact-scales}; the
injection bound is smaller. Subtracting the inviscid equation from
the diffusive one shows that $w^{\mathrm{new}}-w^0$ has zero initial
data and source $\Delta w^0$. Lemma~\ref{lem:torus-scalar} gives
\begin{equation}\label{eq:burst-error}
 \|w^{\mathrm{new}}-w^0\|_\infty
 \le2\delta_nD_n<\frac1{50}<\frac1{16}
\end{equation}
throughout the injection and stirring intervals.

Suppose first that $\Psi_k(p)\in E$ for an integer $k\ge0$.
Choose $n\ge\max\{1,k\}$. At the stirring time with $s_n(t)=k$,
the reference at $Y=\Psi_k(p)$ equals $g_n(p)=1$. The old part is
nonnegative, so~\eqref{eq:burst-error} gives $w(t,Y)>15/16>1/2$.
This proves the first implication of Theorem~\ref{thm:compact-detector}.

Now suppose the orbit stays at distance at least $2d$ from $E$.
Every point of the reference support lies within distance
\[
 \sqrt2\,b_nR^n=\sqrt2\,10^{-6}2^{-n}<1/32\le d
\]
of $\Psi_{s_n(t)}(p)$ during stirring. The same bound follows directly
during injection. The reference support is therefore at distance at
least $d$ from $E$ at every active time. In particular, $w^0=0$
on $E$, and
\begin{equation}\label{eq:active-exclusion}
 w(t,Y)\le HM_0+1/16\qquad(Y\in E)
\end{equation}
during injection and stirring.

At the end of stirring, the reference still has this separation and,
by area preservation, has mass $\int g_n$. During the wait, the new
part is the heat evolution of that reference plus the heat evolution
of an error with supremum norm below $1/16$. The latter bound is
preserved; Lemma~\ref{lem:heat-bounds} bounds the former contribution
on $E$ by $H\int g_n$. Including the old part gives
\begin{equation}\label{eq:wait-exclusion}
 w(t,Y)\le2HM_0+1/16<1/2\qquad(Y\in E).
\end{equation}
The last inequality holds for both pairs in~\eqref{eq:two-constants}.
Before the first block $w=0$. Equations~\eqref{eq:active-exclusion}
and~\eqref{eq:wait-exclusion} cover all later times and every endpoint,
proving the second implication.

\subsection{Effective prescription and physical viscosity}

For a computably supplied $t\ge0$, choose rational $q$ with
$|q-t|\le1$ and put $N_t=\max\{1,\lceil(q+1)/2\rceil\}$.
Every summand of~\eqref{eq:compact-program} with $n>N_t$ starts after
$t$ and vanishes there with all derivatives. The remaining sum is
finite, and its $n$th term contains only $n$ translates of $V_0$.
Effective derivative bounds give $L,C$, and hence all scales.
No comparison of an exact real with a block boundary is needed:
near zero, derivatives of $b$ are polynomial factors times $e^{-1/r}$,
and $s^m e^{-s}\le(m+j)!s^{-j}$ bounds their tails effectively.
Also $b(r)+b(1-r)\ge e^{-2}$. These estimates give effective
evaluation of the smooth profiles even at undecidable endpoint
equalities. Periodic charts are evaluated by finite supersets of their
possibly active translates.

The force formula repeatedly uses the same whole-period processor,
which contains all its local branches. Increasing $n$ asks the flow
to perform more periods, not the force evaluator to determine the
visited configurations. In particular evaluation of $f$ never uses
the unknown scalar $w$.

For a positive computable physical viscosity $\nu$, set
\begin{equation}\label{eq:compact-viscosity}
 U_\nu(t,Y,z)=(\nu a(\nu t,Y),w(\nu t,Y)),\qquad
 f_\nu(t,Y,z)=(\nu^2F(\nu t,Y),\nu h(\nu t,Y)).
\end{equation}
Every horizontal equation acquires the factor $\nu^2$, and every
vertical term the factor $\nu$. Thus direct substitution gives
viscosity $\nu$ with pressure zero. The vertical threshold and mass
are unchanged apart from the time parameter. The time intervals
$\delta_n$ and the waits above are measured in unit-viscosity time;
their physical lengths are divided by $\nu$. Smoothness, effective
evaluation and finite-time uniqueness persist. For an arbitrary
positive real $\nu$, the same formulas have the corresponding
oracle-relative interpretation. This completes the proof of
Theorem~\ref{thm:compact-detector}.

\section{The fixed torus detector}\label{sec:compact-application}

We now supply the geometric input and complete the first part of
Theorem~\ref{thm:main}. The planar processor theorem in
\cite[Theorem~4.1]{stationary2026} effectively constructs from a machine
and its input a smooth one-periodic Hamiltonian field $V$ on $\T^2$,
zero in phase collars, supported in a compact subset of the open unit
square. Its initial point is exactly $p=(1/4,1/4)$. At integer phases
the orbit follows the full closed rectangles of the reversible
compiler. If the machine halts, an integer-phase point enters
\[
 E=\{(x,y):1/32<y<1/8\}
\]
at distance at least $1/32$ from its boundary. If the machine does
not halt, the complete trajectory, including each intermediate
routing segment, satisfies $3/16\le y\le7/8$.
The circle distance from this corridor to $E$ is at least $1/16$.
The same theorem supplies effective global first- and second-derivative
bounds, and its quantitative flow statement supplies both forward and
inverse bounds; alternatively~\eqref{eq:flow-bounds} follows directly
from~\eqref{eq:processor-derivatives}.

Take $d=1/32$ and $H=10^9$. The nonhalting distance is at least $2d$,
and the integer halting entry satisfies the other alternative in
Theorem~\ref{thm:compact-detector}. Its two implications are therefore
exactly the equivalence in~\eqref{eq:torus-event}. All injections are
inside a fixed small square about $(1/4,1/4)$, and the horizontal
force uses only derivatives and products of $V$. Thus their union
is contained in one compact $K\subset(0,1)^2$, fixed for this
prescription. The pressure is zero and the scalar construction gives
the global smooth solution from rest. This proves the first part of
Theorem~\ref{thm:main}.

The second quantitative specialization of the detector theorem is
useful for a planar realization whose nonhalting orbit stays in
$1/8<x<3/8$ and whose target is $2/3<x<5/6$.
The infimum circle separation is at least
\[
 \min\{2/3-3/8,\,1+1/8-5/6\}=7/24>1/4.
\]
Consequently the hypotheses hold with $d=1/16$, $H=10^8$, whenever
that realization also supplies phase collars, integer halting entry
and the derivative bounds in~\eqref{eq:processor-derivatives}.
This is a direct specialization of the proved scalar theorem; it
does not require importing an alternative geometric construction
into the present proof.
This specialization is applied to the planar processor of
\cite[Lemma~7.2]{slowclocks2026} in
\cite[Section~7.5]{slowclocks2026}.
